\begin{otherlanguage*}{english}
These final pages of Chapter~XVIII examine how far accessibility survives constructions which are elementary at the level of categories but delicate at the level of cardinal bounds \citep{Grothendieck1990XVIII}. The passage begins inside the proof of a gluing theorem, moves through categories of sections and categories over a small base, and ends by testing the limits of the theory against subcategories defined by the existence of a morphism. The appended calculation with products of fields supplies the decisive examples. The announced section~4.17 on paratoposes contains no mathematical text in the edition; the reader should not look to that heading for an additional theorem.

One convention governs almost every argument. An indexing category grand devant an infinite cardinal $\pi$ admits cocones for diagrams with at most $\pi$ arrows. Thus the source's $\mathrm L_\pi$ means existence or preservation of $\pi^+$-filtered colimits in the modern regular-cardinal convention. The same successor occurs in a $\pi$-ultrafilter: its intersections may have at most $\pi$ factors. A useful cardinal also controls a small family of presentable objects. These are separate requirements. An accessible category has arbitrarily large suitable ranks, but a proof using a rank chosen in advance must verify the additional uniformity it needs.

The gluing construction explains why these distinctions matter. A category over the walking arrow has an open fibre $\mathcal E_0$, a closed fibre $\mathcal E_1$, and a bifunctor $H(x_0,x_1)$ describing arrows between them. Presentability in a fibre need not be presentability in the glued category: the crossing Hom sets have to preserve the same filtered colimits. The calculations in SGA~4, Expos\'e~I, 9.6 and 9.8, provide the small-limit interchange used in the proof \citep{SGA4I1972}. The audit identifies two necessary repairs. Corollary~4.14.25(1) needs existence of the relevant colimits in the closed fibre. Corollary~4.15.2 needs a common initial rank controlling the stated functors before its subsequent cardinal enlargement. Explicit counterexamples show that these are hypotheses with mathematical content. The existential accessibility theorem survives the repaired choice of ranks.

The categories of sections express the same problem in a more organized form. Products collect the objects, inserters collect the transition maps, and equifiers impose their equations. This relates the source directly to the accessible-category closure results of Makkai and Par\'e, whose book appeared in 1989, and to the contemporary theory of flexible limits. The detailed modern construction in \citet[Appendix~A]{ChingRiehl2014} also explains why this machinery became useful in homotopical algebra. For the total category of a fibration or an opfibration, the commentary gives separate proofs of colimit existence and presentability. The two orientations require different uses of cartesian or cocartesian factorization; merely reversing a statement about filtered colimits would not supply either proof.

The small-base argument is the technical centre of the passage. A sufficiently filtered diagram in a small category can be reduced to one object and one recurrent idempotent. Its formal ind-object is the image of that idempotent, so it becomes representable exactly when the surviving projector splits. The walking unsplit projector, with monoid $\{1,p\}$ and $p^2=p$, is the simplest example. The source eventually proves the strict-bound conclusion
\[
 \operatorname{Ind}_{\pi}(\mathcal C)\simeq
 \operatorname{Kar}(\mathcal C)
 \qquad\bigl(\pi>|\operatorname{Mor}\mathcal C|\bigr).
\]
Here the notation on the left retains the source's successor convention. The commentary supplies both its growth argument and a presheaf proof. The literal identity condition in 4.15.9 is false, even for a two-element group; eventual invertibility and an appropriate trivialization are the invariant conclusions. A separate monoid reduction produces a quotient of a submonoid, not necessarily a quotient of the whole monoid. Neither defect invalidates the strict-bound theorem. The general equality case remains unsettled here, although two-sided cancellation and a stricter local Hom bound give affirmative special cases. Modern comparisons with Cauchy completion and coherent idempotents are provided by \citet{Tendas2022Continuity} and \citet[Sections~5.4.2--5.4.3]{Lurie2009HTT}.

The later questions expose several different ways in which accessibility can fail. The double-dual functor on sets fails to preserve filtered colimits at every rank. Isbell's category of partial bijections is self-dual and finitely accessible, yet has a one-object finite dense and codense subcategory and is not essentially small; this answers all three questions in 4.16.4 negatively \citep[Example~8.7]{AveryLeinster2021}. Other failures depend on the universe. Under the appendix's hypothesis, the commentary proves that the decreasing filtered-colimit closures of an accessible functor's image need not stabilize. A separate argument treats contravariant images. Adjoining retracts of representables to a nonaccessible cosieve also produces a category with a small dense family of presentable objects and sufficiently filtered colimits which is nevertheless not accessible. This repairs the precise defect which Grothendieck notices in his own proposed construction.

The question about existing colimits needs equal care. If all diagrams of a fixed shape have colimits, their colimit cocones form an accessible category, as in \citet[Example~5.4.7.9]{Lurie2009HTT}. The commentary constructs, under the appendix's hypothesis, a finitely accessible category whose existing pushout squares form a nonaccessible full subcategory. It is a collage obtained by adjoining one object to the category of $k$-algebras. Its pushout condition detects exactly the absence of a $k$-point. Thus closure under filtered colimits alone does not produce a small generating family.

The finite tests in 4.16.6 provide a positive result of a different kind. We prove the two split-idempotent tests by first constructing a colimit over the split image and then checking maps out of fibrewise presentable objects. A strict pullback can nevertheless retain an unsplit projector whose image in the base is an identity; the walking isomorphism gives a counterexample to the proposed shortcut b$'$. With all the tests retained, the intended relative property has a concise description. Over a small Karoubian base $\mathcal C$, it is equivalent to accessibility of $\mathcal E$ and $f:\mathcal E\to\mathcal C$, together with the lifting of every prescribed base isomorphism. Such an $f$ is an isofibration. Its strict pullbacks are equivalent to accessible pseudopullbacks, so the proposed change from small test bases to accessible test bases is valid. The standard closure theorem and its applications are discussed in \citet[Examples~4.8 and Remark~4.9, published version]{BourkeLack2024}.

The closing open--closed calculus supplies the conceptual explanation of the counterexamples. Open means a full sieve, closed means a full cosieve. A cosieve is accessible precisely when it has a small weakly initial family, and precisely when its complementary sieve is closed under sufficiently filtered colimits. The canonical slice of presentable objects over $x$ proves these equivalences: a nonempty cosieve in that filtered slice is cofinal, whereas a proper complementary sieve cannot be. This argument repairs the omitted slice denominators in the edition and completes the implication left unfinished in the appendix. The transfinite no-forward-morphism construction is also justified with an explicit universe bound. For the additional question restricted to fields, the answer is exact: in $V_\kappa$ such a $\kappa$-long sequence exists if and only if $\kappa$ is not a Vop\v{e}nka cardinal, with external families allowed. The commentary proves the reduction using a typed graph construction and the fully faithful functor of \citet{MillerPoonenSchoutensShlapentokh2018}.

The appendix identifies every maximal ideal of a product of fields with an ultrafilter. A map from the product back to an infinite ground field $k$ requires $|k|^+$-completeness. Its universe hypothesis says that sufficiently complete nonprincipal ultrafilters do not exist on small sets, equivalently that measurable cardinals below the universe boundary are bounded. Under that hypothesis, for every infinite $\pi\ge |k|$ the algebra
\[
 k^{\pi^+}/\{f:|\operatorname{supp}(f)|\le\pi\}
\]
has no $k$-point, while every subalgebra of size at most $\pi$ has one. A complete proof evaluates representatives outside their common exceptional set. A second construction disproves the proposed bound $\pi=|k|$ for every infinite field without the universe hypothesis. Conversely, a sufficiently large compact cardinal restores a uniform obstruction bound; for finite fields the ordinary ultrafilter lemma already gives a finitely generated obstruction. The comparison with powerful images uses \citet{BrookeTaylorRosicky2017}, without asserting an exact large-cardinal converse for this one fixed-field problem.

The results marked \textsuperscript{[ours]} are deductions developed for this commentary, with full proofs and bounded accounts of the literature checked; the mark is not a priority claim. The remaining limits are stated where they occur, especially the unrestricted cardinal-equality case and accessibility of each individual contravariant closure. All source slips are distinguished from mathematical repairs. This commentary, including its audit and developments, is AI-produced by GPT-6 Astra (OpenAI, ChatGPT Pro mode), checked and integrated by Claude (Anthropic).
\end{otherlanguage*}
